\chapter{Orientation}
\label{ch:orientation}

\section{What problem is being solved}

A generative model is a machine that, having seen a finite collection of examples, produces
new ones that look like they came from the same source. Diffusion models do this by a
two-stage trick: first destroy the data by adding noise gradually until nothing is left but
pure noise, then learn to run that destruction backwards.

Running it backwards requires knowing, at every intermediate noise level, which direction
points back towards plausible data. That direction is called the \emph{score}, and it turns
out to be equivalent to answering a question that sounds much more ordinary:

\begin{quote}
\emph{Given a corrupted example, what was the clean one, on average?}
\end{quote}

That average --- the posterior mean --- is the object everything in this project is about.
A trained diffusion model is, in the end, a machine that estimates it.

\section{The difficulty this project addresses}

For almost every interesting distribution of data, nobody can compute that posterior mean.
There is no formula. So when a neural network is trained to estimate it and makes errors,
there is no way to tell whether the network was bad or the target was hard. The two are
confounded, and that confound blocks a lot of questions one would like to ask about why
diffusion models work.

The few distributions where the posterior mean \emph{is} computable are mostly Gaussians, and
Gaussians are unhelpful here: for them the posterior mean is a linear function of the
observation, so every question about how a model exploits structure has a trivial answer.

This project takes a distribution that is
\begin{itemize}
\item structured enough to be interesting --- the coordinates are correlated;
\item non-Gaussian, so the posterior mean is genuinely nonlinear;
\item and yet exactly computable.
\end{itemize}

The distribution is a \textbf{Markov chain}: a sequence in which each entry depends on the
previous one and, given that, on nothing earlier. Chapter~\ref{ch:markov} defines this
properly. The reason the posterior mean is computable is a structural accident that
Chapter~\ref{ch:bp} explains in full, and which can be previewed in one sentence: noising each
coordinate independently does not create any new dependencies between the clean coordinates,
so the machinery that makes chains tractable survives the noising intact.

\section{The three questions}

\begin{enumerate}
\item \textbf{What does exact inference buy?} If one has the true Markov chain, one can
compute the posterior mean exactly. How much better is that than the standard Gaussian
approximation, and at which noise levels does the difference live?
(Chapters~\ref{ch:bp}, \ref{ch:gaussian}.)

\item \textbf{Can the chain be learned?} In reality the chain is unknown. Can its transition
rule be estimated from noisy observations alone --- never seeing a clean example --- and how
expensive is that? (Later chapters, in progress.)

\item \textbf{How does this compare with a neural network?} Given the same data, does a
structured estimator beat a trained denoiser, and does any advantage in estimating the
posterior mean translate into better generated samples? (Later chapters, in progress.)
\end{enumerate}

The third question has an answer that surprised us, and it is the reason the accompanying note
is framed the way it is. It is stated in the note and derived here once the supporting
machinery is in place.

\section{How to read this document}

Chapters~\ref{ch:probability}--\ref{ch:score} build the diffusion setting from probability
first principles. Chapters~\ref{ch:markov}--\ref{ch:bp} build the graphical-model machinery
and prove the exactness result. Chapter~\ref{ch:gaussian} treats the Gaussian case, which is
both the sanity check and the baseline. Chapter~\ref{ch:numerics} explains what an actual
implementation has to do and where its errors come from.

Three kinds of box appear throughout:

\begin{intuition}
Explains why a result is what it is, in words, before or after the algebra. Skippable if the
algebra was enough.
\end{intuition}

\begin{pitfall}
Names the concrete failure that occurs if an assumption is dropped. These are not
hypothetical --- most of them describe mistakes actually made during this project and caught
later.
\end{pitfall}

\begin{codebox}
Points at the exact file and function implementing the object just defined. Paths are relative
to \codefile{research/nongaussian-bp/} in the repository.
\end{codebox}

\section{Map of the code}
\label{sec:codemap}

Everything referenced in this document lives under \codefile{research/nongaussian-bp/}.

\begin{center}
\small
\begin{tabular}{@{}p{4.4cm}p{10.2cm}@{}}
\toprule
\textbf{Path} & \textbf{Contents} \\
\midrule
\codefile{src/priors.py} & The data-generating models: \texttt{GaussianAR1},
\texttt{LaplaceAR1}, \texttt{StudentTAR1}, \texttt{UniformAR1},
\texttt{GaussianMixtureAR1}. Each provides \texttt{sample} and
\texttt{log\_transition\_matrix}. \\
\codefile{src/bp\_grid.py} & The inference recursion. \texttt{make\_grid},
\texttt{grid\_bp} (one sequence, with evidence), \texttt{grid\_bp\_batch} (many sequences,
device-parameterised). \\
\codefile{src/bp\_gaussian.py} & The closed-form Gaussian recursion in information form. \\
\codefile{src/exact\_scores.py} & Closed-form Gaussian posterior mean and score, used as
ground truth. \\
\codefile{src/em.py} & The expectation step, the sufficient statistic, and the estimation
driver. \\
\codefile{src/kernels.py} & Parametrised transition families and their maximisation steps. \\
\codefile{src/denoiser.py} & A common interface over inference-based and network-based
denoisers, plus the network training loop. \\
\codefile{src/reverse.py} & Reverse-time integration: SDE and probability-flow ODE. \\
\codefile{src/sample\_metrics.py} & Statistics for comparing generated data against the
truth. \\
\codefile{src/backend.py} & CPU/GPU array-module dispatch for the recursion. \\
\codefile{experiments/} & \texttt{exp\_01}--\texttt{exp\_17}, one file per experiment, each
self-describing. \\
\codefile{tests/} & Validation. Several tests here are the evidence for claims made in this
document and are cited as such. \\
\bottomrule
\end{tabular}
\end{center}

\begin{codebox}
Nothing in this document is stated without a pointer to where it is realised. If a derivation
here has no corresponding \textbf{In the code} box, that is deliberate and means the result is
used only on paper.
\end{codebox}
